\begin{remark}[Reproducible: pressure curve \texorpdfstring{$P(q)$}{P(q)}, discrete lower convex hull, and Legendre candidate spectrum]\label{rem:pom-pressure-multifractal-q60}
To convert the pressure curve of Definition \ref{def:pom-moment-pressure} and the subsequent convex conjugate spectral inversion into \emph{runnable numerical objects},
the script \texttt{scripts/exp\_fold\_collision\_pressure\_multifractal\_q60.py} uses the residue DP modulo \(F_{m+2}\) to stably compute \(\log S_q(m)\) on the window \(m\in[24,30]\), and for each integer \(2\le q\le 60\) estimates \(\log r_q\) by linear regression (at \(q\le 17\) the verified exact Perron roots are used as overrides).
The outputs are:
\begin{itemize}
  \item The convergence plot of \(r_q^{1/q}\) (aligned to the endpoint constant \(\sqrt{\varphi}\), Proposition \ref{prop:pom-rq-qinfty-endpoint});
  \item The discrete curve of the pressure \(P(q)=\log r_q-q\log 2\) together with its lower convex hull (consistent with the discrete convexity of Proposition \ref{prop:pom-pressure-convexity}), overlaid with the two linear pinching bounds from Proposition \ref{prop:pom-rq-universal-bounds};
  \item The Legendre candidate spectrum constructed via the convex conjugate
  \[
  f_{\mathrm{Leg}}(\gamma):=\log\varphi-\sup_{q\ge 2}\bigl(q\gamma-\log r_q\bigr),
  \]
  serving as the computable upper bound/candidate curve for the ``spectral inversion'' target described in Proposition \ref{prop:pom-spectrum-legendre-upper} and Remark \ref{rem:pom-spectrum-formalism}.
\end{itemize}
The three figures generated in one run are Figures \ref{fig:fold-collision-rq-1q}--\ref{fig:fold-collision-falpha}.
\input{sections/generated/fig_fold_collision_pressure_multifractal_q60_rq_en}
\input{sections/generated/fig_fold_collision_pressure_multifractal_q60_pressure_en}
\input{sections/generated/fig_fold_collision_pressure_multifractal_q60_falpha_en}
\end{remark}

